% Generated draft: requires mathematical and source review.
\documentclass[11pt]{article}
\usepackage[margin=25mm]{geometry}
\usepackage{fontspec}
\setmainfont{Latin Modern Roman}
\usepackage{amsmath,amssymb,mathtools,unicode-math}
\setmathfont{Latin Modern Math}
\usepackage{enumitem,needspace,xurl,hyperref}
\setlength{\parindent}{0pt}
\setlength{\parskip}{5pt}
\setlength{\emergencystretch}{3em}
\newcommand{\N}{\mathbb{N}}
\newcommand{\Z}{\mathbb{Z}}
\newcommand{\Q}{\mathbb{Q}}
\newcommand{\R}{\mathbb{R}}
\newcommand{\C}{\mathbb{C}}
\newcommand{\F}{\mathbb{F}}
\newcommand{\A}{\mathbb{A}}
\newcommand{\PP}{\mathbb{P}}
\newcommand{\E}{\mathbb{E}}
\newcommand{\eps}{\varepsilon}
\newcommand{\dd}{\,\mathrm{d}}
\begin{document}
\section*{UnsolvedMath ID: 30003906}
\textbf{Source problem code:} catalogue entry
\subsection{Definitions and mathematical statement}
Let \(k\) be a number field, let \(n\geq 1\) be an integer, and let \(G\subseteq S_n\) be a transitive permutation group, where \(S_n\) is the symmetric group on \(\{1,\ldots,n\}\). Transitivity means that for every \(i,j\) there exists \(g\in G\) with \(g(i)=j\). For a real parameter \(X>1\), the supplied record asks for the asymptotic behavior, as \(X\to\infty\), of the schematic count
\[
N_k(G,X)=\#\bigl\{K/k:\operatorname{Gal}(K/k)\simeq G,\ \operatorname{Disc}(K/k)<X\bigr\}.
\]
It further asks whether, in the cases where Malle's original logarithmic exponent is valid,
\[
N_k(G,X)\sim C(G,k)X^{1/a(G)}(\log X)^{b(G,k)-1}.
\]
Here \(\sim\) denotes that the ratio of the two sides tends to \(1\). The expression presupposes a positive constant \(C(G,k)\), independent of \(X\), and specified exponents \(a(G)\neq 0\) and \(b(G,k)\).

\textbf{Unresolved formulation.} These displays preserve the supplied target but do not specify a unique counting problem. The record does not explicitly require \([K:k]=n\), specify whether extensions are counted up to \(k\)-isomorphism, or define the discriminant bound. The relative discriminant of an extension of number fields is an ideal of the ring of integers of \(k\); comparing it with a real number requires a numerical convention, such as its absolute ideal norm, which the record does not state.

The group condition also requires clarification. Requiring \(K/k\) to be Galois with group isomorphic to \(G\) differs from counting degree-\(n\) extensions whose normal closures have Galois group \(G\) in the specified permutation action on the \(n\) embeddings of \(K\). An abstract group isomorphism alone does not encode that permutation action. The record does not resolve these interpretations.

Finally, the record does not define \(a(G)\) or \(b(G,k)\), identify the cases covered by the qualification about the original logarithmic exponent, or address degenerate cases such as the trivial transitive group in degree \(1\). No definitions of these exponents, additional hypotheses, or choice between the counting interpretations are imposed here. The displayed asymptotic is a qualified question attributed to Malle's conjecture, not an asserted theorem.
\subsection{Short English statement}
The record asks how the number of number field extensions with a prescribed permutation group grows with a discriminant bound, and whether a qualified Malle-type asymptotic holds. A precise conjecture requires clarification of the extensions being counted, the discriminant convention, the exponents, and the scope of the logarithmic-exponent qualification. The supplied open label is a dated registry assessment.
\subsection{Sources}
\begin{itemize}
\item Supplied local UnsolvedMath registry record, permanent ID \(30003906\), title \emph{Malle Conjecture for Number Field Counts}, exact source problem code \texttt{catalogue entry}, in the Oberwolfach Reports Open Problems collection. The conjecture is attributed to Malle; no individual proposer is supplied.
\item Supplied report citation: \emph{Explicit Methods in Number Theory (2019)}, Oberwolfach Reports, pp.\ 2033--2076, DOI \texttt{10.4171/owr/2018/34}. Supplied link: \url{https://doi.org/10.4171/owr/2018/34}. The registry separately supplies proposed year \(2018\). The report contents have not been independently consulted for this draft.
\item Supplied literature source: Wang--Lemke Oliver--Wood report, \url{https://ems.press/content/serial-article-files/46754}.
\item Supplied literature source: Alberts--Lemke Oliver--Wang--Wood 2025, \url{https://doi.org/10.1515/crelle-2025-0074}.
\item Literature assessment date: \texttt{2026-08-21}. This is supplied dated curation, not an independently verified assessment. It labels the record open, describes it as an incomplete fragment identifying Malle's conjecture, reports substantial proved families and counterexamples to the original formulation, and concludes that no definite resolution follows for this extracted claim. The supplied background requests independent verification of time-sensitive status and cited claims.
\item The registry labels the statement \texttt{corrected\_verified} and reports that the official Oberwolfach report was checked directly and that definitions were restored. Those are supplied source-verification claims; they have not been independently verified here, and the supplied text still leaves the formulation issues described above unresolved.
\end{itemize}
\end{document}
